\documentclass[12pt,a4paper]{article}

\usepackage[utf8]{inputenc}
\usepackage[T1]{fontenc}


\usepackage{amsmath}
\usepackage{amssymb}
\usepackage{amsthm}
\usepackage[margin=2.5cm]{geometry}
\usepackage{setspace}
\usepackage{natbib}
\usepackage{hyperref}
\usepackage{booktabs}
\usepackage{graphicx}

\hypersetup{
    colorlinks=true,
    linkcolor=black,
    citecolor=black,
    urlcolor=black
}

\onehalfspacing

\title{\textbf{Time--Frequency Methods in Index Tracking}}

\begin{document}

\maketitle

% \begin{abstract}
% The financial econometrics literature has long debated whether correlation
% or cointegration provides the more appropriate foundation for index
% tracking and portfolio replication. We argue that this debate is
% incompletely framed: correlation and cointegration are not competing
% paradigms but limiting cases of a richer spectral structure. By
% recasting the tracking problem in the frequency domain we show that
% cointegration corresponds to a very low-pass constraint on the price
% spread, that correlation-based optimisation targets high-frequency return
% dynamics, and that the two are connected by a continuum parameterised by
% the spectral content of the tracking error. Fractional integration sits
% naturally within this picture, but the frequency-domain representation
% is the cleaner and more practically useful framework, particularly when
% investor horizon preferences are heterogeneous. Time-frequency methods
% from signal processing then allow these preferences to be articulated
% continuously and precisely, and we argue that coherence under coherent
% risk measures is a natural desideratum in this setting. The present
% document establishes the conceptual and mathematical foundations for
% this programme.
% \end{abstract}

% \newpage
% \tableofcontents
% \newpage

%%%%%%%%%%%%%%%%%%%%%%%%%%%%%%%%%%%%%%%%%%%%%%%%%%%%%%%%%%%%%%%%%%%%%%%%%%%
\section{Introduction}
\label{sec:introduction}
%%%%%%%%%%%%%%%%%%%%%%%%%%%%%%%%%%%%%%%%%%%%%%%%%%%%%%%%%%%%%%%%%%%%%%%%%%%

Portfolio optimisation and index replication have long relied on the
statistical relationship between asset prices, yet the financial
econometrics literature has never reached a settled consensus on which
measure of dependence is most appropriate for these tasks. Two candidates
have dominated the debate: correlation, which quantifies the co-movement
of asset returns over a fixed sample window, and cointegration, which
identifies the existence of a common stochastic trend in the price levels
of two or more integrated series. Though both concepts are grounded in
the same probabilistic foundations, they capture fundamentally different
aspects of dependence, and the choice between them has non-trivial
consequences for portfolio construction, risk management, and trading
strategy design.

This paper argues that the debate has been incompletely framed. Flip
flops and mountaineering boots are both footwear. Neither is wrong.
The difficulty arises when one attempts to declare a universal winner
across all terrains: the question of which is better has no answer
until the terrain is specified. The correlation versus cointegration
literature has largely proceeded as though the terrain were fixed,
seeking to settle which approach is superior without first asking
what the investor's horizon actually requires.

The resolution lies in frequency analysis. Correlation and cointegration
are not competing methods but limiting cases of a single family of
tracking criteria, each targeting a different part of the frequency
spectrum of the portfolio-benchmark spread. This becomes precise through
the frequency response of the differencing operator: working in returns
concentrates the tracking objective on fluctuations at timescales of
days, while working in price levels concentrates it on fluctuations at
timescales of years. Neither is universally better. Each is well suited
to a specific investor horizon and poorly suited to all others.

Fractional integration provides a natural interpolation between these two
extremes, with each value of the fractional differencing parameter $d$
corresponding to a concrete intermediate horizon. But the fractional
family is itself limited: it can only produce monotone frequency
responses, and real investor preferences are not always monotone. A
general frequency response relaxes this restriction entirely, allowing
the tracking criterion to be shaped to any horizon profile. The result
is a framework in which investor time horizon preference is a design
parameter, expressible quantitatively and continuously, rather than a
binary choice between two inherited methodologies. The present work
proposes such a framework, grounded in time-frequency representations
of the price process. This is extended to linear risk-coherent
optimisation programmes \citep{ArtznerDelbaenEberHeath1999} for both index tracking and enhanced index
tracking, and portfolio sparsity is enforced through the $\ell_0$-norm
approximation procedure of \citet{BenidisFengPalomar2018}.

\subsection{The Correlation Paradigm and Its Limitations}

The correlation-based paradigm has its roots in the mean-variance
framework of \citet{Markowitz1959} and was formalised for index tracking
by \citet{Roll1992}, whose tracking error variance minimisation model
remains the practitioner standard. The appeal of correlation is its
familiarity and computational tractability: given a sample of returns, a
covariance matrix can be estimated and an optimal portfolio derived in
closed form. The limitation, however, is equally well understood.
Correlation is computed from return series, which are stationary by
construction, and therefore captures only short-run co-movements.
\citet{Alexander1999} observed that correlation is intrinsically a
short-run measure and that tracking portfolios built on it carry no
guarantee that the price spread between the replica and the benchmark
will remain bounded over time. For daily data, this short-run character
is pronounced: the tracking objective is primarily sensitive to
fluctuations at horizons of days to weeks, and is essentially blind to
slower-moving divergences between the portfolio and the benchmark. This
follows mathematically from working with return series, which are first
differences of prices: differencing acts as a high-pass filter,
suppressing low-frequency trends and retaining only rapid fluctuations
\citep{ChaudhuriLo2016,EngleGranger1974}.
Section~\ref{sec:frequency_analysis}
makes this precise and quantifies the effective horizon of
correlation-based tracking. In the absence of a mean-reversion
mechanism, the tracking error may behave as a random walk, drifting
arbitrarily far from zero and necessitating frequent and costly
rebalancing. This observation, which has since been reproduced in
numerous empirical settings, motivates the cointegration alternative.

\subsection{The Cointegration Framework}

The cointegration framework, introduced by \citet{Granger1981} and
\citet{EngleGranger1987} and later extended by \citet{Johansen1988,
Johansen1991}, addresses the deficiency of correlation by working
directly with price levels. Two or more $I(1)$ series are said to be
cointegrated if a stationary linear combination of them exists. In the
index tracking context this is a natural property: because a market index
is itself a linear combination of its constituent prices, there must in
principle exist a portfolio whose value tracks the index with a
stationary, mean-reverting error. \citet{Alexander1999} formalised this
insight and showed that a cointegration-optimal portfolio guarantees
stationarity of the tracking error by construction, ties the portfolio to
its benchmark in the long run, and produces more stable portfolio weights
because they are estimated from price levels over a long calibration
window. For daily data, calibration windows of three to five years are
standard, and the method is largely insensitive to fluctuations at
horizons shorter than a year. This long-run character is the structural
counterpart to correlation's short-run sensitivity: working in price
levels, which are the cumulative sum of returns, acts as a low-pass
filter, suppressing rapid fluctuations and retaining only slow-moving
common trends \citep{ChaudhuriLo2016}. Section~\ref{sec:frequency_analysis} quantifies this
precisely. A direct consequence is that cointegration-based portfolios
require less frequent rebalancing than their correlation-based
counterparts. The cause is structural: the objective is itself
insensitive to the short-run deviations that would otherwise trigger
rebalancing.

These theoretical properties were subsequently developed into explicit
trading strategies by \citet{AlexanderGiblinWeddington2002} and
\citet{AlexanderDimitriu2002, AlexanderDimitriu2005a,
AlexanderDimitriu2005b}, who demonstrated that cointegration-based
long-short equity strategies can achieve low volatility, near
market-neutrality, and favourable Sharpe ratios across multiple market
regimes. The error correction model that underpins cointegration reveals
causal flows, sometimes called Granger causalities, that must exist
between series sharing a common stochastic trend \citep{Granger1988}, and
this dynamic structure is itself a source of useful information for
portfolio management and risk control.

\subsection{An Open and Inconclusive Empirical Debate}
\label{sec:litreview}

\subsubsection{Evidence from the Literature}

The empirical evidence, however, is considerably less clear-cut than the
theoretical arguments would suggest. \citet{AlexanderDimitriu2005a}, in
what remains the most rigorous direct comparison, found that for simple
index tracking using Dow Jones Industrial Average constituents the
out-of-sample performance of cointegration-optimal portfolios was broadly
similar to that of tracking error variance minimisation, with neither
method exhibiting a decisive advantage once the number of stocks was
sufficient. \citet{Grobys2010}, applying the same comparison to the
Swedish OMX market over a ten-year period, reached the opposite
conclusion, finding that cointegration-based portfolios clearly dominated
on both Sharpe and Treynor ratios, with the best portfolio outperforming
the index by a substantial margin at lower volatility. \citet{SantAnna2016},
who made portfolio selection endogenous to the cointegration optimisation
through large-scale simulation across Brazilian and United States data,
found again that neither method consistently dominates: cointegration
delivered better tracking error for Brazilian equities while correlation
was superior on turnover and volatility, and results for the United
States market were mixed.

In a different application of the same two building blocks,
\citet{Miao2014} proposed a two-stage procedure in which correlation
filters candidate pairs and cointegration provides the final selection
criterion for high-frequency statistical arbitrage, reporting strong
out-of-sample performance on energy sector stocks. The combination of
the two stages is telling: it suggests that each measure captures
something the other misses, and that a unified treatment would be
preferable to sequential filtering. \citet{Arsov2023}, examining small
South-East European markets, found that although Johansen tests
initially indicated cointegration, the result could not be validated
diagnostically, and the analysis reverted to vector autoregressive and
GARCH frameworks. The broader literature on emerging market integration
and volatility spillovers similarly finds that the presence or absence
of cointegration is sensitive to market structure, data frequency, and
sample length, with no universal verdict.

Beyond the index tracking literature, similar ambiguity appears in
studies of international market linkages. A recurring finding is that
developed markets transmit both return and volatility shocks to smaller
markets through mechanisms that operate at multiple time scales
simultaneously \citep{Arsov2023, AlexanderDimitriu2005b}. These
multi-scale transmission patterns cannot be adequately characterised by
either correlation, which aggregates across all frequencies, or
cointegration, which focuses exclusively on the lowest.

\subsubsection{A Question of Investor Horizon}

The inconclusiveness of this literature is less surprising once the
horizon argument is taken seriously. The studies above largely evaluate
both methods against the same short-run tracking error metrics,
regardless of the horizon the method was designed for. This is precisely
where the terrain matters. Asking which method is better without
specifying the investor's horizon is not a well-posed question.

Consider two investors. The first is a portfolio manager running a
short-term equity overlay, rebalancing weekly to maintain index exposure
for hedging purposes. Daily and weekly deviations from the benchmark
directly affect the hedge's effectiveness and need to be kept small.
Multi-year drift is irrelevant: positions are rarely held beyond a few
months. For this investor, correlation-based tracking is the natural
choice. The tracking objective is sensitive to exactly the fluctuations
that matter, and frequent rebalancing is part of how the strategy
operates rather than a cost to be avoided.

The second is an asset manager running the equity sleeve of a defined
benefit pension fund, matching assets to liabilities that extend twenty
years into the future. Whether the portfolio was 0.3\% away from the
benchmark on any given Tuesday is of no consequence. What matters is
that when liabilities fall due, the portfolio value is close to the
benchmark value. Rebalancing costs compound over decades, so stability
of weights is a primary concern. For this investor, cointegration-based
tracking is the natural choice. The long-run adherence it guarantees is
exactly what the liability-matching mandate requires, and the lower
turnover is a direct benefit.

Neither investor is wrong about their method. They are simply operating
at different points on the frequency spectrum of the tracking error.
The empirical debate, by evaluating both methods on a common metric
without regard to horizon, has been asking which of the two serves both
investors simultaneously, and correctly finding that neither does so
consistently.

This reframing points naturally toward an intermediate territory. A
quarterly-reporting fund, a retail investor with a five-year savings
horizon, or a risk-parity manager rebalancing monthly all sit between
these two extremes and are served by neither method in its pure form.
The theory of fractional integration, introduced by
\citet{GrangerJoyeux1980} and \citet{Hosking1981}, provides a
continuous interpolation. By allowing the order of differencing $d$ to
take non-integer values between zero and one, it produces tracking
criteria whose effective horizon varies continuously from days to years
as $d$ decreases from one toward zero. Correlation corresponds to
$d = 1$; cointegration corresponds to $d \approx 0$; and every
investment horizon in between corresponds to some intermediate value.
Section~\ref{sec:fractional_integration} develops this connection
formally, and Section~\ref{sec:frequency_analysis} makes the horizon
mapping precise.

\subsection{A False Dichotomy}

What emerges from this body of work is an open and incompletely specified
debate. The studies reviewed above have largely asked
which method is better, without first specifying better for whom and over
what horizon. As the two investor examples in the previous section
illustrate, this question has no horizon-independent answer: correlation
is the correct tool for the equity overlay manager hedging weekly
exposures, and cointegration is the correct tool for the pension fund
matching twenty-year liabilities. Evaluated on the same short-run
tracking error metric, cointegration will appear inferior. Evaluated on
long-run price-level adherence, correlation will appear inferior. Neither
conclusion is informative about the methods themselves; both reflect the
mismatch between the evaluation criterion and the method's intended
horizon. This is the primary source of the inconclusiveness documented
in Section~\ref{sec:litreview}.

Beyond this misspecification, the discussion has also proceeded in binary
terms: a portfolio is either built on correlation or on cointegration,
the tracking error is either stationary or it is not, and the analyst
either differences the data or regresses in levels. Consider a third
investor: a target-date retirement fund manager with a ten-year horizon,
rebalancing quarterly and reporting annual performance to clients. Daily
noise is irrelevant to this investor, but so is the exclusive focus on
decade-scale price-level adherence that cointegration enforces. This
investor sits squarely between the two extremes, and the binary framing
offers no natural home for them. The dichotomy is therefore a second and
separate limitation, compounding the primary one.

Correlation and cointegration are two limiting cases of a single, richer
structure. Correlation, by
operating on returns, is sensitive to high-frequency co-movements while
being blind to low-frequency common trends. Cointegration, by identifying
stochastic trends in price levels, is precisely a statement about the
very lowest frequencies of the data generating process: the near-zero
portion of the spectrum where power accumulates without bound in an
integrated series. The two approaches therefore sample different parts of
the spectral content of the multivariate price process.

This observation connects naturally to the theory of fractional
integration. A series that is $I(0)$ is stationary and its spectral
density is bounded and flat at the origin. A series that is $I(1)$ has a
spectral density that diverges at zero frequency, reflecting the
accumulation of permanent shocks. Fractional integration, in the sense
of \citet{GrangerJoyeux1980} and \citet{Hosking1981}, allows the order
of integration to take non-integer values between zero and one, producing
series whose spectral density diverges at the origin at a rate that
interpolates continuously between the stationary and the unit-root cases.
From this perspective, the correlation-versus-cointegration debate is a
debate about whether the tracking error lives at $d = 0$ or at some
positive fractional value approaching unity, and the practitioner's
choice of optimisation method implicitly encodes an assumption about
which frequency bands are relevant to the investment horizon in question.
Section~\ref{sec:fractional_integration} makes this mapping precise,
providing a concrete correspondence between values of $d$ and investment
horizons in trading days.

The fractional differencing family, however, imposes a fundamental
restriction: its frequency responses are necessarily monotone. For any
value of $d$, the resulting filter either up-weights low frequencies or
high frequencies, and cannot selectively target an intermediate band.
An investor whose mandate is defined at a specific horizon, such as a
quarterly-reporting fund that is indifferent to both daily noise and
multi-year drift, cannot express this preference through any choice of
$d$. The generalisation to arbitrary frequency responses, developed in
Section~\ref{sec:general_freq}, removes this restriction
entirely.

\subsection{Structure of the Paper}
\label{sec:structure}

The remainder of this paper is organised as follows.
Section~\ref{sec:foundations} establishes the mathematical foundations
required for the analysis. It introduces fractional differentiation and
its relationship to long-memory processes
(Section~\ref{sec:fractional_integration}), presents the discrete
Fourier transform and the frequency response as analytical tools
(Section~\ref{sec:fourier}), and motivates the move to time-frequency
representations (Section~\ref{sec:tfr}). It then sets out the index
tracking and enhanced index tracking optimisation problems
(Section~\ref{sec:it_formulation}), reviews coherent risk measures and
their linear programme formulations (Section~\ref{sec:risk_measures}),
and introduces sparse index tracking (Section~\ref{sec:sparsity}).

Section~\ref{sec:freq_tracking} develops the core argument
quantitatively, mirroring the narrative of Section~\ref{sec:introduction}
with precise derivations. It derives the frequency responses of the
differencing and integration operators and connects them to concrete
investor horizons via the standard signal processing half-power cutoff
(Sections~\ref{sec:corr_response} and~\ref{sec:coint_response}). It
then derives the fractional differencing family and maps each value of
$d$ to a trading-day horizon (Section~\ref{sec:frac_response}). It
demonstrates that passband representations allow more precise horizon
targeting than the monotone fractional family permits
(Section~\ref{sec:passband}), and argues that the non-stationarity of
financial returns necessitates time-frequency rather than purely
Fourier-based analysis (Section~\ref{sec:stationarity}).

%%%%%%%%%%%%%%%%%%%%%%%%%%%%%%%%%%%%%%%%%%%%%%%%%%%%%%%%%%%%%%%%%%%%%%%%%%%
\section{Foundations}
\label{sec:foundations}
%%%%%%%%%%%%%%%%%%%%%%%%%%%%%%%%%%%%%%%%%%%%%%%%%%%%%%%%%%%%%%%%%%%%%%%%%%%

\subsection{Fractional Differentiation}
\label{sec:fractional_integration}

The distinction between the correlation and cointegration paradigms rests
on whether one works with returns or with prices, that is, on whether the
price series is differenced once or not at all. This is a binary choice
between two integer orders of differencing. Fractional differentiation
generalises the differencing operation to non-integer orders, and in
doing so replaces the binary choice with a continuous one.

It is convenient to express differencing through the lag operator $L$,
defined by its action of shifting a series one step into the past,
\begin{equation}
    L x_t = x_{t-1}, \qquad L^k x_t = x_{t-k}.
\end{equation}
The first difference of a series, the change from one period to the next,
is then
\begin{equation}
    \Delta x_t = x_t - x_{t-1} = (1 - L) x_t,
\end{equation}
so that the operation of differencing is represented by the operator
$(1 - L)$. Differencing $d$ times corresponds to applying this operator
$d$ times, that is, to raising it to the power $d$.

Because $(1-L)^d$ is a binomial raised to a power, it expands according to
the binomial theorem,
\begin{equation}
    (1 - L)^d = \sum_{k=0}^{d} \binom{d}{k} (-L)^k,
\end{equation}
where the coefficients $\binom{d}{k}$ are the familiar binomial
coefficients. For $d = 2$, for example, the expansion yields $(1-L)^2 = 1
- 2L + L^2$, so that $(1-L)^2 x_t = x_t - 2x_{t-1} + x_{t-2}$. For any
positive integer $d$ the sum terminates at $k = d$, since the binomial
coefficients vanish beyond that point, and differencing therefore
involves only a finite number of past observations.

The binomial coefficient is conventionally defined through factorials,
\begin{equation}
    \binom{d}{k} = \frac{d!}{k!\,(d-k)!}.
\end{equation}
This definition restricts $d$ to non-negative integers, since the
factorial is defined only for such values. To admit non-integer orders of
differencing, a generalisation of the factorial to arbitrary real
arguments is required.

Such a generalisation is provided by the gamma function, introduced by
\citet{Euler1738}, which extends the factorial to the real and complex
numbers while preserving its defining property,
\begin{equation}
    \Gamma(n + 1) = n! \quad \text{for non-negative integers } n,
\end{equation}
and remaining well defined for non-integer arguments. Substituting the
gamma function for the factorials in the binomial coefficient gives
\begin{equation}
    \binom{d}{k} = \frac{\Gamma(d + 1)}{\Gamma(k + 1)\,\Gamma(d - k + 1)},
\end{equation}
an expression that coincides with the ordinary binomial coefficient when
$d$ is an integer but remains meaningful for any real $d$.

With the binomial coefficient defined for arbitrary real orders, the
difference operator extends directly to non-integer $d$. Following
\citet{GrangerJoyeux1980} and \citet{Hosking1981}, the fractional
difference operator of order $d$ is defined by the series
\begin{equation}
    (1 - L)^d = \sum_{k=0}^{\infty} \binom{d}{k} (-L)^k
    = \sum_{k=0}^{\infty}
    \frac{\Gamma(d + 1)}{\Gamma(k + 1)\,\Gamma(d - k + 1)} (-1)^k L^k.
\end{equation}
Unlike the integer case, the series no longer terminates: for non-integer
$d$ the coefficients are never exactly zero, and the operator depends on
the entire past of the series. Applied to a series, fractional
differentiation is thus a weighted sum of all preceding values,
\begin{equation}
    (1 - L)^d x_t = x_t - d\, x_{t-1}
    + \frac{d(d-1)}{2}\, x_{t-2} - \cdots,
\end{equation}
with weights determined by the generalised binomial coefficients.

The dependence on the entire history raises a question of practical
computability, since only a finite sample is ever available. The weights,
however, decay as $k$ increases, so the operator can be approximated to
arbitrary accuracy by truncating the series at a finite number of lags.

Figure~\ref{fig:fracdiff_sp500} illustrates the effect of the operator on
the S\&P 500 price series for several orders of differencing. At $d = 0$
the series is the raw price level; at $d = 1$ it is the ordinary return
series, stationary and dominated by short-term fluctuation. The
intermediate orders $d = 0.25$ and $d = 0.75$ produce series that retain
progressively less of the slow-moving trend as $d$ increases, tracing a
continuous path between the two integer extremes.

\begin{figure}[h]
\centering
\includegraphics[width=\textwidth]{figures/fracdiff_sp500.pdf}
\caption{The S\&P 500 price series under fractional differencing of
orders $d = 0$ (the raw price level), $d = 0.25$, $d = 0.75$, and $d = 1$
(ordinary returns). As $d$ increases from zero to one, the slow-moving
trend is progressively removed and the series becomes increasingly
dominated by short-term fluctuation.}
\label{fig:fracdiff_sp500}
\end{figure}

\subsection{Fourier Analysis and Frequency Responses}
\label{sec:fourier}

\subsubsection{The Discrete Fourier Transform}
\label{sec:dft}

Let $x_0, x_1, \ldots, x_{T-1}$ be a finite sequence of real-valued
observations, for instance the daily returns of an asset over $T$ trading
days. The discrete Fourier transform re-expresses this sequence as a
combination of complex exponentials, each oscillating at a distinct
frequency. The transform is defined by
\begin{equation}
    X_k = \sum_{t=0}^{T-1} x_t \, e^{-j 2\pi k t / T},
    \qquad k = 0, 1, \ldots, T-1,
\end{equation}
where $j = \sqrt{-1}$ and each index $k$ corresponds to the angular
frequency $\omega_k = 2\pi k / T$ in radians per sample. The coefficient
$X_k$ is a complex number whose magnitude $|X_k|$ measures the amplitude
of the oscillation at frequency $\omega_k$ and whose argument $\angle X_k$
measures its phase, that is, the position of the oscillation relative to
the start of the sample. A large $|X_k|$ indicates that a cyclical
component with period $T/k$ samples accounts for a substantial part of
the variation in the series.

The original sequence is recovered without loss through the inverse
transform,
\begin{equation}
    x_t = \frac{1}{T} \sum_{k=0}^{T-1} X_k \, e^{j 2\pi k t / T},
    \qquad t = 0, 1, \ldots, T-1.
\end{equation}
Writing the complex exponential through Euler's identity, $e^{j\theta} =
\cos\theta + j\sin\theta$, and using the conjugate symmetry $X_{T-k} =
\overline{X_k}$ that holds because the input is real-valued, the inverse
transform can be expressed entirely in terms of real trigonometric
functions,
\begin{equation}
    x_t = \frac{1}{T} \sum_{k=0}^{T-1}
    \Big( \Re[X_k]\cos\tfrac{2\pi k t}{T}
    - \Im[X_k]\sin\tfrac{2\pi k t}{T} \Big),
\end{equation}
or equivalently, collecting the magnitude and phase of each coefficient,
as
\begin{equation}
    x_t = \frac{1}{T} \sum_{k=0}^{T-1}
    |X_k| \cos\!\Big( \tfrac{2\pi k t}{T} + \angle X_k \Big).
\end{equation}
This last form makes the interpretation explicit: the series is a sum of
cosine waves, one for each frequency $\omega_k$, with amplitude
$|X_k|/T$ and phase offset $\angle X_k$. The decomposition is exact
rather than approximate, so no information is lost in moving between the
time representation $\{x_t\}$ and the frequency representation $\{X_k\}$.
The lowest non-zero frequency, $k = 1$, corresponds to a single cycle
spanning the entire sample of $T$ observations, while the highest, $k =
\lfloor T/2 \rfloor$, corresponds to an oscillation that alternates at
nearly every observation.

\subsubsection{The Power Spectrum and the Decomposition of Variance}
\label{sec:power_spectrum}

The Fourier coefficients also reveal how the variability of a series is
distributed across timescales. The link is a classical result known as
Parseval's theorem, which states that the total energy of a signal is the
same whether measured in the time domain or the frequency domain. For a
zero-mean series this energy is exactly the sample variance, and the
theorem takes the form
\begin{equation}
    \operatorname{var}(x_t) = \frac{1}{T} \sum_{t=0}^{T-1} x_t^2
    = \frac{1}{T^2} \sum_{k=1}^{T-1} |X_k|^2.
\end{equation}
The practical content of this identity is that the variance of the series
can be split into separate, additive contributions, one from each
frequency. Define the periodogram
\begin{equation}
    \hat{S}_{xx}[k] = \frac{1}{T^2} |X_k|^2,
\end{equation}
which is simply the share of the total variance carried by the
oscillation at frequency $\omega_k$. Summing these shares over all
frequencies returns the variance in full; summing them over a subset of
frequencies returns the part of the variance that lives in that band. The
periodogram is therefore a budget: it tells the analyst exactly where the
movement in a series comes from.

This is the heart of the matter, and it can be stated without any
equation at all. Given a return series, one may ask what fraction of its
total variation arises from rapid fluctuations between consecutive
observations and what fraction arises from slow movements unfolding over
many months. The power spectrum answers precisely this question. It takes
the total variance, the single quantity into which ordinary risk
measurement collapses all variation, and distributes it across a range of
timescales, showing how much of an asset's variability is attributable to
fast dynamics and how much to slow.

A few concrete shapes make the idea tangible. If a series were pure white
noise, with no structure linking one observation to the next, its power
would be spread evenly across all frequencies: the spectrum would be
flat, and no timescale would carry more variance than any other. This is
the benchmark against which real data is judged. A series dominated by
slow trends, by contrast, concentrates its power at the low-frequency
end: most of its variance comes from movements that take months or years
to unfold, while the day-to-day variation is comparatively small. A
series dominated by rapid reversals shows the opposite, with power piled
up at the high-frequency end. The location of power along the frequency
axis is thus a fingerprint of the dynamics of the series.

The translation from frequency to calendar time is direct and is what
makes the spectrum economically meaningful. A frequency $\omega_k$
corresponds to a cycle that repeats every $T/k$ observations, so for
daily data a given point on the frequency axis maps to a definite number
of trading days. Power concentrated at cycles of a few days reflects the
rapid, largely transitory fluctuations associated with market
microstructure and short-term reversals. Power at cycles of weeks to
months reflects the medium-term dynamics associated with fund flows,
rebalancing, and momentum. Power at cycles of a year or more reflects the
slow components associated with business cycles, monetary regimes, and
secular trends. When an analyst inspects the power spectrum of a return
series, they are reading off how the asset's risk is apportioned among
these horizons, an apportionment that the single summary number of
variance conceals entirely.

This decomposition is the reason the frequency domain is the natural
setting for the questions raised in Section~\ref{sec:introduction}. Two
assets, or an asset and a benchmark, may carry the same total variance and
yet distribute it across timescales in entirely different ways. An
investor who cares about one horizon and not another is, in effect,
interested in one region of the spectrum and indifferent to the rest. To
make that preference precise, the variance must first be unpacked into its
frequency components, which is exactly what the power spectrum provides.
The next step is to extend the same decomposition from the variance of a
single series to the covariance between two, since it is covariance with
the benchmark that index tracking ultimately seeks to control.

\subsubsection{The Cross-Spectrum and the Decomposition of Covariance}
\label{sec:cross_spectrum}

The same logic that splits the variance of one series across frequencies
splits the covariance between two series in exactly the same way. This is
the step that matters for index tracking, because tracking is
fundamentally a statement about covariance: a portfolio replicates a
benchmark when the two move together, and the quality of replication is
governed by how their co-movement is distributed across timescales.

Let $x_t$ and $y_t$ be two zero-mean series, for instance the returns of
an asset and the returns of a benchmark, with Fourier coefficients $X_k$
and $Y_k$. The sample covariance decomposes as
\begin{equation}
    \operatorname{cov}(x_t, y_t) = \frac{1}{T} \sum_{t=0}^{T-1} x_t y_t
    = \frac{1}{T^2} \sum_{k=1}^{T-1} \Re[X_k \overline{Y_k}].
\end{equation}
Each term in this sum is the contribution that frequency $\omega_k$ makes
to the total covariance. Define the cross-periodogram
\begin{equation}
    \hat{S}_{xy}[k] = \frac{1}{T^2} \Re[X_k \overline{Y_k}],
\end{equation}
so that the covariance is the sum of these frequency-specific
contributions, and the covariance restricted to a band of frequencies
$\mathcal{K}$ is simply the sum taken over that band. Just as the
periodogram is a budget for variance, the cross-periodogram is a budget
for covariance: it shows at which timescales two series move together and
at which they do not.

The crucial difference from the variance case is that the contributions
can be negative. A frequency where the two series rise and fall together,
in phase, contributes positively to the covariance. A frequency where one
rises while the other falls, out of phase, contributes negatively. The
total covariance is the net result of summing these signed contributions
across all timescales. This is why two series can have low overall
covariance for entirely different reasons: either they barely move
together at any frequency, or they move together strongly at some
frequencies and oppositely at others, with the contributions cancelling
in the sum.

Consider an asset that tracks a benchmark closely on a day-to-day basis
but slowly drifts away from it over the course of years, perhaps because
it belongs to a sector in secular decline. Its high-frequency co-movement
with the benchmark is strongly positive: at the scale of days, the two
are almost indistinguishable. Its low-frequency co-movement, however, is
weak or even negative: at the scale of years, the asset and the benchmark
part ways. The cross-periodogram of this pair would show large positive
values at high frequencies and small or negative values at low
frequencies. Now consider a second asset with the opposite profile: noisy
and only loosely coupled to the benchmark from one day to the next, but
firmly anchored to it over the long run, so that the two never wander far
apart in level. This pair shows the reverse pattern, with weak
high-frequency co-movement and strong positive low-frequency co-movement.

These two assets pose entirely different propositions to an index
tracker, and which one is preferable depends entirely on the investor's
horizon. The first asset is an excellent daily tracker and a poor
long-term one; the second is a poor daily tracker and an excellent
long-term one. An ordinary covariance calculation, which sums across all
frequencies and reports a single number, cannot distinguish between them
in this way: it might even assign them similar total covariances with the
benchmark while concealing that the covariance is generated at opposite
ends of the spectrum. The cross-spectrum makes the distinction explicit,
and in doing so makes it possible to select assets according to the
timescale at which co-movement with the benchmark is actually required.

This is the formal counterpart to the correlation and cointegration
discussion of Section~\ref{sec:introduction}. The covariance of returns,
on which the correlation paradigm is built, weights all frequencies of
co-movement together, but because returns are differenced prices, that
weighting is heavily tilted toward the high-frequency end, as the next
section will make precise. Cointegration, by contrast, is a statement
about co-movement at the very lowest frequencies, where the price levels
of two integrated series remain tied together. The two paradigms are
therefore reading different entries in the same covariance budget.

\subsubsection{Linear Filters and the Reshaping of the Spectrum}
\label{sec:filters}

The decompositions established in the previous two sections characterise
the frequency content of a series as given. It remains to introduce the
mechanism by which that content may be modified, so that specific
timescales can be emphasised and others suppressed. This is the function
of the linear filter.

A linear time-invariant filter transforms an input series $x_t$ into an
output series $y_t$ through a weighted sum of past and present values,
\begin{equation}
    y_t = \sum_{\tau} h_\tau \, x_{t-\tau},
\end{equation}
where the sequence $\{h_\tau\}$ is the impulse response that
characterises the filter. In the time domain this operation, a
convolution, combines neighbouring observations and its effect is not
readily apparent. In the frequency domain it admits a simple description.
The convolution theorem states that filtering multiplies each Fourier
coefficient by a complex factor that depends only on the frequency,
\begin{equation}
    Y_k = H(\omega_k) \, X_k, \qquad
    H(\omega) = \sum_{\tau} h_\tau \, e^{-j\omega\tau},
\end{equation}
where $H(\omega)$ is the frequency response of the filter. The defining
property is that the filter does not transfer energy between frequencies;
it scales each frequency independently. The oscillation at frequency
$\omega_k$ in the output is the corresponding oscillation in the input,
multiplied by the single factor $H(\omega_k)$.

The implication for the variance and covariance budgets is direct. Since
each coefficient is scaled by $H(\omega_k)$, the periodogram of the
output equals the periodogram of the input scaled by $|H(\omega_k)|^2$,
and the cross-periodogram between two filtered series is scaled
identically. Filtering therefore reweights the frequency-by-frequency
contributions to variance and covariance, attenuating those timescales at
which $|H(\omega)|$ is small and preserving those at which it approaches
unity. The profile of $|H(\omega)|^2$ specifies exactly which timescales
the filter retains and which it removes, so that the selection of a
filter is equivalent to the selection of a weighting of the covariance
budget.

Two filters anchor the analysis that follows, and they occupy opposite
ends of the spectrum. The first is the moving average, which replaces
each observation by the mean of the surrounding $L$ values. Its frequency
response is
\begin{equation}
    H_{\text{MA}}(\omega) = \frac{1}{L} \sum_{\tau=0}^{L-1} e^{-j\omega\tau},
    \qquad
    |H_{\text{MA}}(\omega)| = \frac{1}{L}\left|
    \frac{\sin(\omega L / 2)}{\sin(\omega / 2)} \right|,
\end{equation}
which equals unity at zero frequency and decays toward zero as frequency
increases, with the transition governed by the window length $L$. The
moving average is accordingly a low-pass filter, passing slow movements
and suppressing rapid ones, which formalises the sense in which smoothing
removes short-term fluctuations. It follows that the covariance of two
series each smoothed by a moving average equals the sum of the
cross-spectrum over only the low frequencies the filter preserves.
Smoothing the data prior to measuring covariance constitutes
frequency-selective covariance estimation concentrated at long
timescales.

The second filter is the first difference, which replaces each
observation by its change from the preceding period, $y_t = x_t -
x_{t-1}$. Its frequency response is
\begin{equation}
    H_\Delta(\omega) = 1 - e^{-j\omega}, \qquad
    |H_\Delta(\omega)| = 2\left|\sin\tfrac{\omega}{2}\right|,
\end{equation}
which vanishes at zero frequency and increases monotonically to its
maximum at the highest frequency. Differencing is therefore a high-pass
filter, removing slow movements and amplifying rapid ones, in direct
contrast to the moving average. Measuring covariance on differenced data
concentrates the estimate at short timescales.
Figure~\ref{fig:filter_responses} plots the magnitude responses of the
two filters on a common axis, making the low-pass and high-pass
characters and their mirrored structure apparent.

\begin{figure}[h]
\centering
\includegraphics[width=0.8\textwidth]{figures/filter_responses.pdf}
\caption{Magnitude responses of the moving-average (low-pass) and
first-difference (high-pass) filters as functions of angular frequency.
The moving average preserves low frequencies and attenuates high ones,
while the first difference does the reverse; the two are mirror images
across the frequency band.}
\label{fig:filter_responses}
\end{figure}

These two filters are the operations that distinguish the two paradigms
of Section~\ref{sec:introduction}.
Correlation-based tracking operates on returns, which are differenced
prices, and therefore measures co-movement through a high-pass filter
that emphasises the shortest timescales. Cointegration-based tracking
operates on price levels, which are the cumulative sum of returns and
hence the output of the extreme low-pass filter represented by
integration, and therefore measures co-movement at the lowest
frequencies. The two paradigms constitute the same covariance budget
viewed through opposite filters. Section~\ref{sec:freq_tracking} derives
these frequency responses in full, establishes the timescales to which
they correspond under daily sampling, and demonstrates that the range of
filters between them yields a continuum of tracking criteria, each
matched to a distinct investor horizon.

\subsection{Time-Frequency Representations}
\label{sec:tfr}

\subsection{Index Tracking and Enhanced Index Tracking}
\label{sec:it_formulation}

\subsection{Coherent Risk Measures and Linear Programme Formulations}
\label{sec:risk_measures}

\subsection{Sparse Index Tracking}
\label{sec:sparsity}

%%%%%%%%%%%%%%%%%%%%%%%%%%%%%%%%%%%%%%%%%%%%%%%%%%%%%%%%%%%%%%%%%%%%%%%%%%%
\section{Frequency-Aware Index Tracking}
\label{sec:freq_tracking}
%%%%%%%%%%%%%%%%%%%%%%%%%%%%%%%%%%%%%%%%%%%%%%%%%%%%%%%%%%%%%%%%%%%%%%%%%%%

\subsection{Frequency Response of the Differencing Operator}
\label{sec:corr_response}

\subsection{Frequency Response of the Integration Operator}
\label{sec:coint_response}

\subsection{The Fractional Differencing Family}
\label{sec:frac_response}

\subsection{Passband Representations and Investor Horizon Targeting}
\label{sec:passband}
\label{sec:general_freq}

\subsection{Stationarity and the Need for Time-Frequency Representations}
\label{sec:stationarity}

\bibliographystyle{plainnat}

\begin{thebibliography}{99}

\bibitem[Artzner et al.(1999)]{ArtznerDelbaenEberHeath1999}
Artzner, P., Delbaen, F., Eber, J.-M., and Heath, D. (1999).
\newblock Coherent measures of risk.
\newblock \emph{Mathematical Finance}, 9(3), 203--228.

\bibitem[Benidis et al.(2018)]{BenidisFengPalomar2018}
Benidis, K., Feng, Y., and Palomar, D.~P. (2018).
\newblock Sparse portfolios for high-dimensional financial index tracking.
\newblock \emph{IEEE Transactions on Signal Processing}, 66(1), 155--170.

\bibitem[Alexander(1999)]{Alexander1999}
Alexander, C. (1999).
\newblock Optimal hedging using cointegration.
\newblock \emph{Philosophical Transactions of the Royal Society of London,
  Series A}, 357(1758), 2039--2058.

\bibitem[Chaudhuri and Lo(2016)]{ChaudhuriLo2016}
Chaudhuri, S.~E., and Lo, A.~W. (2016).
\newblock Spectral portfolio theory.
\newblock \emph{Working Paper}, MIT Laboratory for Financial Engineering.

\bibitem[Engle(1974)]{EngleGranger1974}
Engle, R.~F. (1974).
\newblock Band spectrum regression.
\newblock \emph{International Economic Review}, 15(1), 1--11.

\bibitem[Euler(1738)]{Euler1738}
Euler, L. (1738).
\newblock De progressionibus transcendentibus seu quarum termini generales
  algebraice dari nequeunt.
\newblock \emph{Commentarii Academiae Scientiarum Petropolitanae}, 5,
  36--57.

\bibitem[Gencay et al.(2002)]{GencaySelcukWhitcher2002}
Gen\c{c}ay, R., Sel\c{c}uk, F., and Whitcher, B. (2002).
\newblock \emph{An Introduction to Wavelets and Other Filtering Methods in
  Finance and Economics}.
\newblock Academic Press, San Diego.

\bibitem[Alexander and Dimitriu(2002)]{AlexanderDimitriu2002}
Alexander, C., and Dimitriu, A. (2002).
\newblock The cointegration alpha: Enhanced index tracking and long-short
  equity market neutral strategies.
\newblock \emph{ISMA Centre Discussion Papers in Finance}, 2002-08.
\newblock University of Reading.

\bibitem[Alexander and Dimitriu(2005a)]{AlexanderDimitriu2005a}
Alexander, C., and Dimitriu, A. (2005a).
\newblock Indexing and statistical arbitrage: Tracking error or cointegration?
\newblock \emph{The Journal of Portfolio Management}, 31(2), 50--63.

\bibitem[Alexander and Dimitriu(2005b)]{AlexanderDimitriu2005b}
Alexander, C., and Dimitriu, A. (2005b).
\newblock Indexing, cointegration and equity market regimes.
\newblock \emph{International Journal of Finance and Economics}, 10(3),
  213--231.

\bibitem[Alexander et al.(2002)]{AlexanderGiblinWeddington2002}
Alexander, C., Giblin, I., and Weddington, W. (2002).
\newblock Cointegration and asset allocation: A new active hedge fund strategy.
\newblock \emph{Research in International Business and Finance}, 16, 65--90.

\bibitem[Arsov(2023)]{Arsov2023}
Arsov, S. (2023).
\newblock Cointegration across stock markets in South-East Europe and possible
  spillovers.
\newblock \emph{SSRN Working Paper}. Available at SSRN~4640959.

\bibitem[Engle and Granger(1987)]{EngleGranger1987}
Engle, R.~F., and Granger, C.~W.~J. (1987).
\newblock Co-integration and error correction: Representation, estimation, and
  testing.
\newblock \emph{Econometrica}, 55(2), 251--276.

\bibitem[Granger(1981)]{Granger1981}
Granger, C.~W.~J. (1981).
\newblock Some properties of time series data and their use in econometric
  model specification.
\newblock \emph{Journal of Econometrics}, 16(1), 121--130.

\bibitem[Granger(1988)]{Granger1988}
Granger, C.~W.~J. (1988).
\newblock Some recent developments on a concept of causality.
\newblock \emph{Journal of Econometrics}, 39(1--2), 199--211.

\bibitem[Granger and Joyeux(1980)]{GrangerJoyeux1980}
Granger, C.~W.~J., and Joyeux, R. (1980).
\newblock An introduction to long memory time series models and fractional
  differencing.
\newblock \emph{Journal of Time Series Analysis}, 1(1), 15--29.

\bibitem[Grobys(2010)]{Grobys2010}
Grobys, K. (2010).
\newblock Correlation versus cointegration: Do cointegration based index
  tracking portfolios perform better? Evidence from the Swedish stock market.
\newblock \emph{German Journal for Young Researchers}, 2(1), 72--78.

\bibitem[Hosking(1981)]{Hosking1981}
Hosking, J.~R.~M. (1981).
\newblock Fractional differencing.
\newblock \emph{Biometrika}, 68(1), 165--176.

\bibitem[Johansen(1988)]{Johansen1988}
Johansen, S. (1988).
\newblock Statistical analysis of cointegration vectors.
\newblock \emph{Journal of Economic Dynamics and Control}, 12(2--3), 231--254.

\bibitem[Johansen(1991)]{Johansen1991}
Johansen, S. (1991).
\newblock Estimation and hypothesis testing of cointegration vectors in
  Gaussian vector autoregressive models.
\newblock \emph{Econometrica}, 59(6), 1551--1581.

\bibitem[Markowitz(1959)]{Markowitz1959}
Markowitz, H. (1959).
\newblock \emph{Portfolio Selection: Efficient Diversification of Investments}.
\newblock John Wiley and Sons, New York.

\bibitem[Miao(2014)]{Miao2014}
Miao, G.~J. (2014).
\newblock High frequency and dynamic pairs trading based on statistical
  arbitrage using a two-stage correlation and cointegration approach.
\newblock \emph{International Journal of Economics and Finance}, 6(3), 96--110.

\bibitem[Roll(1992)]{Roll1992}
Roll, R. (1992).
\newblock A mean/variance analysis of tracking error.
\newblock \emph{The Journal of Portfolio Management}, 18(4), 13--22.

\bibitem[Sant'Anna et al.(2016)]{SantAnna2016}
Sant'Anna, L.~R., Filomena, T.~P., and Caldeira, J.~F. (2016).
\newblock Index tracking and enhanced indexing using cointegration and
  correlation with endogenous portfolio selection.
\newblock \emph{The Quarterly Review of Economics and Finance}, 65, 146--157.

\end{thebibliography}

\end{document}